% Level 3, Week 5 -- Sorting, Stacks and Queues (student workbook)
\chapterweek{5}{Sorting, Stacks and Queues}{Sorting puts data in order; a stack remembers what to undo; a queue makes sure everyone gets a turn.}

\section{What you will learn}
\begin{checklist}
  \item Sort a list with bubble sort and selection sort, and count the swaps.
  \item Sort a leaderboard with \texttt{sorted(key=..., reverse=True)}.
  \item Use a stack to undo, and a queue to take turns.
\end{checklist}

\section{Two ways to sort}
\textbf{Selection sort:} find the smallest item and swap it to the front. Then find the smallest of the rest and swap it into second place. Repeat.\\
\textbf{Bubble sort:} walk along the list comparing neighbors. If the left one is bigger, swap them. After each pass the biggest item has \enquote{bubbled} to the end.

\begin{center}
%% alt: Sorting 5, 2, 4, 1 in passes. Start: 5 2 4 1. After pass 1: 2 4 1 5, with 5 in its final place. After pass 2: 2 1 4 5, with 4 in place. After pass 3: 1 2 4 5, sorted. Green boxes with a thick border are in their final place.
\begin{tikzpicture}[card/.style={draw=black,line width=0.8pt,rounded corners=2pt,minimum width=0.8cm,minimum height=0.9cm,font=\small\bfseries}]
  \foreach \row/\a/\b/\c/\d/\note/\done in {
      0/5/2/4/1/{start}/0,
      1/2/4/1/5/{after pass 1: 5 is in place}/1,
      2/2/1/4/5/{after pass 2: 4 is in place}/2,
      3/1/2/4/5/{after pass 3: sorted!}/4}{
    \foreach \v [count=\k from 0] in {\a,\b,\c,\d}{
      \pgfmathtruncatemacro{\placed}{\k >= 4-\done ? 1 : 0}
      \ifnum\placed=1
        \node[card,fill=kcgreen!35,line width=1.8pt] at (\k*1.0,-1.15*\row) {\v};
      \else
        \node[card,fill=white] at (\k*1.0,-1.15*\row) {\v};
      \fi
    }
    \node[anchor=west,font=\small] at (3.8,-1.15*\row) {\note};
  }
  \node[anchor=west,font=\small] at (3.8,0.8) {green, thick border = in its final place};
\end{tikzpicture}
\end{center}

\begin{workedexample}[Selection sort by hand]
Sort the cards 5, 2, 4, 1 with selection sort. Each pass looks at the part that is not sorted yet.
\begin{center}
\renewcommand{\arraystretch}{1.25}
\begin{tabular}{@{}llll@{}}
\toprule
\textbf{Pass} & \textbf{Smallest of the rest} & \textbf{What happens} & \textbf{Row after the pass} \\ \midrule
1 & 1 & swap 1 with 5 (1 swap) & \textbf{1}, 2, 4, 5 \\
2 & 2 & already in place, no swap & \textbf{1, 2}, 4, 5 \\
3 & 4 & already in place, no swap & \textbf{1, 2, 4, 5} \\
\bottomrule
\end{tabular}
\end{center}
Selection sort needed only \textbf{1} swap. Bubble sort needed \textbf{5} swaps for the same cards (look at the picture above: 3 in pass 1, 1 in pass 2, 1 in pass 3). Same answer, different amount of work.
\end{workedexample}

\begin{tip}[Your turn: sort 4, 1, 3, 2 both ways]
Use 4 paper cards. Write the row after each pass and count the swaps.\\[2pt]
Selection sort: pass 1 \rule{2.2cm}{0.4pt}\ pass 2 \rule{2.2cm}{0.4pt}\ pass 3 \rule{2.2cm}{0.4pt}\quad swaps \rule{1cm}{0.4pt}\\[4pt]
Bubble sort:\ \ \ pass 1 \rule{2.2cm}{0.4pt}\ pass 2 \rule{2.2cm}{0.4pt}\ pass 3 \rule{2.2cm}{0.4pt}\quad swaps \rule{1cm}{0.4pt}
\end{tip}

\section{Python's own sort}
Python can sort for you. The \textbf{key} is a function that says \emph{what to sort by}:
\begin{pycode}
scores = [("Ada", 120), ("Leo", 95), ("Mei", 150)]

def points_of(entry):
    return entry[1]                  # sort by the number

best_first = sorted(scores, key=points_of, reverse=True)
print(best_first)   # [('Mei', 150), ('Ada', 120), ('Leo', 95)]
\end{pycode}
\texttt{sorted} calls \texttt{points\_of} on every entry: \texttt{("Ada", 120)} gives 120, \texttt{("Leo", 95)} gives 95, \texttt{("Mei", 150)} gives 150. Then it puts the entries in order of those numbers. \texttt{reverse=True} means \enquote{biggest first}. Write \texttt{key=points\_of} with \textbf{no} parentheses: you hand Python the function, and Python calls it for you.
\texttt{sorted} makes a \textbf{new} list and leaves the old one alone.

\section{Stacks and queues}
\begin{center}
%% alt: A stack and a queue. A stack holds stroke 1, stroke 2 and stroke 3 (top); you push and pop at the top, so it is last in, first out. A queue holds Priya at the front, then Sam, Kai and Zoe; new people join at the back (append) and leave from the front (popleft), so it is first in, first out.
\begin{tikzpicture}[item/.style={draw=black,line width=0.8pt,minimum width=2.2cm,minimum height=0.6cm,font=\small},
                    q/.style={draw=black,line width=0.8pt,rounded corners=4pt,minimum width=1.15cm,minimum height=0.8cm,font=\small}]
  % stack
  \node[item,fill=kcsky!25] at (0,0) {stroke 1};
  \node[item,fill=kcsky!25] at (0,0.6) {stroke 2};
  \node[item,fill=kcorange!40,line width=1.4pt] at (0,1.2) {stroke 3 (top)};
  \draw[line width=1pt] (-1.25,1.6) -- (-1.25,-0.35) -- (1.25,-0.35) -- (1.25,1.6);
  \draw[-{Stealth[length=3mm]},line width=1pt] (2.0,2.3) -- node[right,font=\small]{push} (0.8,1.65);
  \draw[-{Stealth[length=3mm]},line width=1pt] (-0.8,1.65) -- node[left,font=\small]{pop} (-2.0,2.3);
  \node[font=\small\bfseries] at (0,-0.9) {Stack: last in, first out};
  % queue
  \foreach \n/\x in {Priya/6, Sam/7.3, Kai/8.6, Zoe/9.9}{\node[q,fill=kcyellow!40] at (\x,0.6) {\n};}
  \draw[-{Stealth[length=3mm]},line width=1pt] (5.4,0.6) -- (4.4,0.6) node[left,font=\small,align=center]{leaves\\(popleft)};
  \draw[-{Stealth[length=3mm]},line width=1pt] (11.6,0.6) -- (10.6,0.6);
  \node[font=\small,align=center] at (11.6,1.35) {joins\\(append)};
  \node[font=\small] at (5.9,1.35) {front};
  \node[font=\small\bfseries] at (7.8,-0.9) {Queue: first in, first out};
\end{tikzpicture}
\end{center}
A \textbf{stack} is like a pile of plates: you can only take the top one. \enquote{Undo} uses a stack: the last thing you did is the first thing undone. A \textbf{queue} is like a lunch line: join at the back, leave from the front. In Python, \texttt{collections.deque} is a fast queue.
\begin{pycode}[A stack and a queue in Python]
from collections import deque

plates = []                # a list works as a stack
plates.append("red")       # push: put on top
plates.append("blue")
print(plates.pop())        # pop: take the top -> blue

line = deque()             # a deque works as a queue
line.append("Priya")       # join at the back
line.append("Sam")
print(line.popleft())      # leave from the front -> Priya
\end{pycode}

\begin{exercise}[Sort Race]
Folder: \texttt{week05/ex-sort-race}. \texttt{selection\_sort} is done. Finish the swap inside \texttt{bubble\_sort}:
\begin{pycode}
if result[i] > result[i + 1]:
    # swap both at once:
    result[i], result[i + 1] = result[i + 1], result[i]
    swaps = swaps + 1
\end{pycode}
Run it. The program sorts the 8 cards \texttt{[7, 3, 9, 1, 5, 3, 8, 2]} both ways. Write the swaps it reports: selection \rule{1.5cm}{0.4pt}\quad bubble \rule{1.5cm}{0.4pt}\\
\textbf{Done looks like:} both lines show \texttt{[1, 2, 3, 3, 5, 7, 8, 9]} (both 3s are still there), and \texttt{check\_my\_work.py} passes.
\end{exercise}

\begin{exercise}[Leaderboard]
Folder: \texttt{week05/ex-leaderboard}. Finish \texttt{name\_of} (sort A to Z, ignoring capitals), \texttt{top\_scores} and the fair tie rule in \texttt{rank\_lines}: equal points share a place, so the board reads 1, 2, 3, 3, 5. It works like a race: if two runners tie for 3rd place, nobody is 4th, and the next runner is 5th.\\
\textbf{Done looks like:} the full board shows \texttt{3. Leo - 95} and \texttt{3. Omar - 95}, then \texttt{5. Zoe - 80}.
\end{exercise}

\begin{exercise}[Undo Paint]
Folder: \texttt{week05/ex-undo-paint}. Finish the \texttt{Stack} class: \texttt{push}, \texttt{pop}, \texttt{peek} and \texttt{size}. Then run it, click the window, draw with the arrow keys and press \texttt{u} to undo. Each stroke is one item on the stack.\\
\textbf{Done looks like:} each \texttt{u} removes the newest line only, the stroke count at the top goes down by 1, and pressing \texttt{u} on an empty picture does nothing (no crash).
\end{exercise}

\begin{tip}[Kind leaderboards]
A good leaderboard uses nicknames, lets ties share a place, and celebrates personal bests and \enquote{most improved}, not only the winner.
\end{tip}

\begin{vocab}
  \vocabword{sort}{Put items in order, for example smallest to largest.}
  \vocabword{swap}{Two items change places.}
  \vocabword{key function}{A function that tells \texttt{sorted} what to sort by.}
  \vocabword{stack}{A collection where the last item in is the first out (push and pop).}
  \vocabword{queue}{A collection where the first item in is the first out (join at the back, leave at the front).}
  \vocabword{deque}{Python's fast queue, from \texttt{collections}. Say \enquote{deck}.}
\end{vocab}

\begin{reflection}
Where do you see a stack and a queue in real life or in apps you use?
\answerlines{2}
Why is \enquote{undo} a stack and not a queue?
\answerlines{1}
\end{reflection}
